\section{The two experiments and their information difference}
\label{sec:r-experiments}

The comparison begins with two experiments that give the signal and
the finite message the same joint distribution.  They differ in the
rows revealed alongside that message.  Before defining them, we fix
the measure and entropy conventions needed to interpret their exact
labels.

Write \(\sigma\) for uniform probability on \(S^{d-1}\), put \(n=d-1\),
and let \(B(z,t)\) be the closed Euclidean ball with center \(z\) and
radius \(t\).  The notation \(\gamma_a\) denotes the law of an
\(a\times d\) matrix with independent \(N(0,1)\) entries.  All
logarithms are natural.  Constants denoted by \(C\) are positive and
absolute and may increase between occurrences.

\begin{lemma}[Spherical caps and ambient balls]\label{lem:r-caps}
On the unit sphere in \(\R^a\), \(a\ge2\), a cap of angular radius
\(0\le\theta\le\pi\) has uniform probability at most
\(\theta^{a-1}\).  On \(S^{d-1}\),
\[
 \sigma(B(s,t))\le t^n\quad(s\in S^{d-1}),\qquad
 \sigma(B(z,t))\le(2t)^n\quad(z\in\R^d),\qquad t>0.
\]
Consequently a fixed unit-vector output, or a random unit-vector
output independent of the signal, succeeds to angular error
\(\epsilon\) on a set of \(\sigma\)-mass at most \(\epsilon^n\).
\end{lemma}

\begin{proof}
The cap probability is the ratio of the integrals of
\(\sin^{a-2}u\) over \([0,\theta]\) and \([0,\pi]\).  The numerator
is at most \(\theta^{a-1}/(a-1)\).  Symmetry and
\(\sin u\ge2u/\pi\) on \([0,\pi/2]\) bound the denominator below by
\(\pi/(a-1)\), proving the angular estimate.

For the centered Euclidean estimate it suffices to consider \(t\le1\).
Rotate \(s\) to the north pole.  Points in \(B(s,t)\) have last
coordinate at least \(1-t^2/2\ge1/2\).  As a graph above the first
\(n\) coordinates, this part of the sphere has area element at most
twice Lebesgue measure and projects into the \(n\)-ball of radius
\(t\).  The entire sphere has area at least twice the volume of the
unit \(n\)-ball, by projecting the two hemispheres.  Their ratio is
at most \(t^n\).  For \(t>1\), total mass one suffices.  An ambient
ball meeting the sphere is contained, on the sphere, in the ball of
radius \(2t\) about any one point of that intersection.  This gives
the second bound.  Finally, angular error at most \(\epsilon\)
implies Euclidean error at most \(\epsilon\), since
\(2\sin(\theta/2)\le\theta\).  Integrating the centered bound over
an independent output law proves the last assertion.
\end{proof}

For probability measures \(P,Q\), relative entropy means
\[
 D(P\Vert Q)=
 \begin{cases}
 \displaystyle\int\log\frac{dP}{dQ}\,dP,&P\ll Q,\\
 +\infty,&\text{otherwise}.
 \end{cases}
\]
For a countable discrete variable \(W\), \(H(W)\) is its Shannon
entropy in nats.  Mutual information and conditional mutual
information are the corresponding relative entropies of a joint law
from the product of its marginals, conditionally when appropriate.
We use regular conditional laws on standard Borel spaces.  The
conditional chain rule and average conditional relative-entropy
formulas used here are given by Austin
\cite[Lemma~3.6 and Equations~(13)--(18)]{Austin2020}.
For probability data processing under a measurable map or a common
Markov kernel, see van Erven and Harremo\"es
\cite[Theorems~5 and~9; Example~2]{vanErvenHarremoes2014}.
Our finite-reference comparisons use the following precise
normalization.

\begin{lemma}[A finite reference measure]\label{lem:r-finite-reference}
Let \(P\) be a probability measure and let \(Q\) be a finite
nonzero measure on the same measurable space.  Put \(c=Q(\Omega)\)
and \(\overline Q=Q/c\).  When \(P\ll Q\), define
\[
 D_{\mathrm{fin}}(P\Vert Q)=\int\log\frac{dP}{dQ}\,dP.
\]
This extended quantity is bounded below by \(-\log c\) and satisfies
\[
 D_{\mathrm{fin}}(P\Vert Q)
     =D(P\Vert\overline Q)-\log c.
\]
If \(T\) is measurable, then
\[
 D_{\mathrm{fin}}(T_\#P\Vert T_\#Q)
 \le D_{\mathrm{fin}}(P\Vert Q).
\]
\end{lemma}

\begin{proof}
The density identity \(dP/dQ=c^{-1}dP/d\overline Q\) gives the
displayed equality and lower bound.  Both \(Q\) and \(T_\#Q\) have
mass \(c\), and \(\overline{T_\#Q}=T_\#\overline Q\).
Probability data processing applied to \((P,\overline Q)\), followed
by subtracting the same \(\log c\) from both sides, proves the claim.
For completeness, probability data processing follows from the
identity
\[
 \frac{dT_\#P}{dT_\#\overline Q}(T)
   =\E_{\overline Q}\left[\frac{dP}{d\overline Q}\,\middle|\,T\right]
\]
and conditional Jensen for \(x\log x\).  Thus the finite comparison
uses a normalized probability reference before the inequality is
applied.
\end{proof}

We will also use the entropy variational inequality.  If \(P\ll Q\)
are probabilities and \(f\) is bounded and measurable, then
\begin{equation}\label{eq:r-entropy-variational}
 \E_P f\le D(P\Vert Q)+\log\E_Q e^f.
\end{equation}
Indeed, normalize \(e^fQ\) to a probability measure and expand the
nonnegative relative entropy from \(P\) to that measure.  Truncation
gives the same statement for nonnegative \(f\) whenever
\(\E_Qe^f<\infty\), and for integrable \(f\) with a finite
exponential moment.  This is the standard variational relation
between entropy and exponential moments; see
Dupuis and Mao \cite[Equations~(1.1)--(1.2), pp.~1--2]{DupuisMao2022}.

\subsection{Density versions for exact labels}
\label{subsec:r-densities}

A Gaussian block maps a spherical signal into Euclidean label space.
The next lemma specifies a common measurable density version and the
integrability that permits subsequent entropy subtractions.  Its
first-moment argument is also a separate route from cube averages to
the exact density.

\begin{lemma}[Projection densities and logarithmic integrability]
\label{lem:r-projection-density}
Let \(1\le\ell<n\).  If \(\rho\) is a finite measure with
\(\rho\ll\sigma\), then the image \(g_\#\rho\) of \(\rho\) under
every full-row-rank \(g\in\R^{\ell\times d}\) has a Lebesgue density.
For \(\delta>0\), put
\[
 p_{\rho,\delta}(g,y)
   =(2\delta)^{-\ell}\int
       \ind_{\{\norm{gu-y}_\infty\le\delta\}}\,d\rho(u).
\]
A jointly Borel density version is
\begin{equation}\label{eq:r-density-version}
 p_\rho(g,y)=\lim_{j\to\infty}p_{\rho,1/j}(g,y)
\end{equation}
where the limit exists and is finite, and zero otherwise.

Suppose now that \(\nu\) is a probability measure with
\(\nu\le H\sigma\), \(H<\infty\), and \(G\sim\gamma_\ell\) is
independent of \(S\sim\nu\).  Then \(\log p_\nu(G,GS)\) is
integrable.  In addition,
\begin{equation}\label{eq:r-first-density-moment}
 \sup_{\delta>0}\E p_{\nu,\delta}(G,GS)
 \le(2\pi)^{-\ell/2}
       \iint\norm{u-s}^{-\ell}\,d\nu(u)d\nu(s)<\infty,
\end{equation}
and the same upper bound holds for \(\E p_\nu(G,GS)\).

Finally, let a countable \(W\) with \(H(W)<\infty\) have any joint
law with \((S,G)\) that leaves the latter's product marginal
unchanged.  For each positive-probability \(w\), put
\(\nu_w=P_{S\mid W=w}\).  The actual conditional law of \(GS\)
given \((G,W)\) has a Lebesgue density.  The logarithm of that
density and \(\log p_{\nu_W}(G,GS)\) are both integrable under the
actual law.
\end{lemma}

\begin{proof}
The first \(\ell\) coordinates of a uniform spherical point have
density
\[
 \frac{\Gamma(d/2)}
      {\pi^{\ell/2}\Gamma((d-\ell)/2)}
 (1-\norm{y}^2)^{(d-\ell-2)/2}
 \ind_{\{\norm{y}<1\}}.
\]
To see the normalization, write the point as a standard Gaussian
vector divided by its length.  Its projected squared length has the
\(\operatorname{Beta}(\ell/2,(d-\ell)/2)\) law and its projected
direction is uniform.  Polar coordinates give the displayed
density.  An orthogonal change of coordinates followed by the
invertible map \((gg^{\mathsf T})^{1/2}\) gives the density for
\(g_\#\sigma\), with Jacobian
\(\det(gg^{\mathsf T})^{-1/2}\).  Absolute continuity gives a
density for \(g_\#\rho\).  Each cube average is Borel in \((g,y)\)
by integration of a Borel indicator.  Lebesgue differentiation
shows that \eqref{eq:r-density-version} is a density version for
each full-rank \(g\); the specified limit is jointly Borel.

We first verify logarithmic integrability from this explicit
spherical density.  Under the product law \(\gamma_\ell\otimes\sigma\),
the logarithm of the residual squared radius is integrable by the
beta law.  Gram--Schmidt expresses \(\det(GG^{\mathsf T})\) as a
product of chi-square variables with positive degrees of freedom,
whose logarithms have finite absolute means.  Hence
\(\log p_\sigma(G,GS)\) is integrable.  Domination by \(H\)
preserves this integrability under \(\gamma_\ell\otimes\nu\).
For every full-rank \(g\), data processing gives
\[
 D(g_\#\nu\Vert g_\#\sigma)\le D(\nu\Vert\sigma)\le\log H.
\]
The negative logarithmic part of any probability density ratio
\(a/b\), integrated under \(a\), is at most \(1/e\), because
\(t\log^-(t)\le1/e\).  Its positive part is therefore finite when
the relative entropy is finite.  Apply this to \(p_\nu/p_\sigma\)
and add the integrable baseline logarithm.

Here is a direct proof that retains the first density moment in
\eqref{eq:r-first-density-moment}.  For distinct \(u,s\), the
\(\ell\) coordinates of \(G(u-s)\) are independent centered
normals with variance \(\norm{u-s}^2\).  Their joint density is
at most \((2\pi)^{-\ell/2}\norm{u-s}^{-\ell}\).  The cube volume
is \((2\delta)^\ell\).  Multiplying these quantities and applying
Tonelli proves the first inequality in
\eqref{eq:r-first-density-moment}; the diagonal is null since
\(\nu\le H\sigma\) is nonatomic.  For each \(s\in S^{d-1}\),
Lemma~\ref{lem:r-caps} and dyadic shells give
\[
 \int\norm{u-s}^{-\ell}\,d\nu(u)
 \le 1+H2^\ell\sum_{j\ge0}2^{-j(n-\ell)}<\infty.
\]
The term \(1\) covers distances greater than one, and \(n-\ell>0\)
makes the series converge.  This is uniform in \(s\).

For each full-rank \(G\), Lebesgue differentiation holds at \(GS\)
for \(\nu\)-almost every \(S\), since \(G_\#\nu\) has a density.
Fubini and Fatou pass the first-moment bound to \(p_\nu(G,GS)\).
It follows that \(\E\log^+p_\nu(G,GS)<\infty\).  To control the
other part, let \(\phi_\ell\) be the standard Gaussian density on
\(\R^\ell\).  For each \(G\),
\[
 \int p_\nu(G,y)
       \left(\log\frac{p_\nu(G,y)}{\phi_\ell(y)}\right)_-dy
 \le\frac1e.
\]
Also \(\E\norm{GS}^2=\ell\), so
\(\E[-\log\phi_\ell(GS)]=\frac{\ell}{2}\log(2\pi)+\frac{\ell}{2}\).
The inequality
\[
 (-\log p_\nu)_+
 \le\left(\log(p_\nu/\phi_\ell)\right)_-
        +(-\log\phi_\ell)_+
\]
proves the negative logarithmic bound directly.  The densities in
these conclusions describe the exact labels; the cube averages only
prove their estimates.

For the last assertion, conditioning on \(W=w\) dominates the joint
law of \((G,S)\) by \(P(W=w)^{-1}\) times
\(\gamma_\ell\otimes\nu\).  This gives an actual conditional image
density for every full-rank \(G\), up to the usual null set in \(G\).
The product marginal is unchanged after averaging over \(w\), so the
actual expectation of \(\lvert\log p_\nu(G,GS)\rvert\) is finite.
The relative entropy of the actual conditional image density from
\(p_\nu(G,\cdot)\), averaged over \((G,W)\), equals
\(I(W;GS\mid G)\le H(W)\).  The density-ratio negative part is at
most \(1/e\), hence its logarithm is integrable and so is the actual
log density.

Similarly, conditional data processing gives
\[
 \E D\bigl(P_{GS\mid G,W}\Vert G_\#\nu_W\bigr)
 \le \E D\bigl(P_{S\mid G,W}\Vert\nu_W\bigr)
 =I(S;G\mid W)\le H(W).
\]
The last inequality follows from \(I(S;G)=0\) and the chain rule.
The log ratio in this display is integrable by the same negative-part
bound.  Subtracting it from the actual log density proves integrability
of \(\log p_{\nu_W}(G,GS)\).  The argument also applies to a
deterministic refinement \(W'=(W,E(S,W))\) whenever \(H(W')<\infty\).
\end{proof}

\subsection{A message selected using actual rows}
\label{subsec:r-two-experiments}

Let \(m,r\ge1\), \(\ell=m+r<n\), and let \(S\sim\sigma\).
Let \(A\sim\gamma_m\) be independent of \(S\), and let \(W\) be a
countable message with \(H(W)<\infty\), coupled arbitrarily to
\((S,A)\).  We define two laws:

\begin{itemize}
\item In the \emph{aligned experiment} \(P_1\), retain the joint law
of \((S,A,W)\), take \(C\sim\gamma_r\) independently, and set
\(G=(A;C)\) and \(Y=GS\).
\item In the \emph{independent experiment} \(P_0\), retain the
marginal law of \((S,W)\), draw \(G\sim\gamma_\ell\) independently
of that pair, and again set \(Y=GS\).
\end{itemize}

Both laws have the same \((S,W)\) marginal and the same product
\((S,G)\) marginal.  In the aligned law, \(W\) may be dependent on
the rows \(G\); in the independent law, \(G\) is independent of
\((S,W)\).  Define
\[
 J_i(W)=I_i(S;W\mid G,Y),\qquad i=0,1,
\]
and, for positive-probability \(w\), let
\(\nu_w=P_{S\mid W=w}\), which is common to both experiments.

\begin{proposition}[The information identity and its density bound]
\label{prop:r-information-comparison}
In the two experiments just defined, all quantities below are finite,
and
\begin{align}
 J_0(W)-J_1(W)
 &=h_0(Y\mid G,W)-h_1(Y\mid G,W)-I_1(S;G\mid W)
 \label{eq:r-information-identity}\\
 &\le h_0(Y\mid G,W)+\E_1\log p_{\nu_W}(G,Y).
 \label{eq:r-information-density-bound}
\end{align}
If \(E=E(S,W)\) is countable and \(H(E\mid W)<\infty\), and the
same deterministic refinement \(W'=(W,E)\) is used in both laws,
then
\begin{equation}\label{eq:r-refinement-charge}
 J_0(W)-J_1(W)
 \le J_0(W')-J_1(W')+H(E\mid W).
\end{equation}
\end{proposition}

\begin{proof}
Because \(Y\) is determined by \((S,G)\), the conditional chain
rule gives
\[
 J_i(W)=I_i(S;W\mid G)-I_i(W;Y\mid G).
\]
Since \(S\) and \(G\) are independent in both laws,
\[
 I_i(S;W\mid G)=I(S;W)+I_i(S;G\mid W),
\]
and \(I_0(S;G\mid W)=0\).  The \((S,W)\) and \((S,G)\)
marginals agree, so subtraction gives
\[
 J_0(W)-J_1(W)
 =I_1(W;Y\mid G)-I_0(W;Y\mid G)-I_1(S;G\mid W).
\]
All mutual information terms are finite: those involving the
countable message are at most \(H(W)\), and the last is at most
\(H(W)\) by this same chain rule.  Lemma~\ref{lem:r-projection-density}
makes the conditional differential entropies finite.  Expressing
the two \(I(W;Y\mid G)\) terms as entropy differences, whose
\(h(Y\mid G)\) terms agree, proves
\eqref{eq:r-information-identity}.  In particular the conditional
row-information term has a negative sign.

Apply conditional data processing to \(s\mapsto Gs\):
\begin{align*}
 -h_1(Y\mid G,W)-\E_1\log p_{\nu_W}(G,Y)
 &=\E_1 D(P_1(Y\in\cdot\mid G,W)\Vert G_\#\nu_W)\\
 &\le\E_1 D(P_1(S\in\cdot\mid G,W)\Vert\nu_W)\\
 &=I_1(S;G\mid W).
\end{align*}
The integrability assertion of the same lemma justifies these
subtractions.  Substituting this inequality into
\eqref{eq:r-information-identity} proves
\eqref{eq:r-information-density-bound}.

Finally
\[
 J_i(W')=J_i(W)+I_i(S;E\mid W,G,Y).
\]
The difference for \(W\) therefore equals the difference for \(W'\)
minus the \(i=0\) term and plus the \(i=1\) term.  The former is
nonnegative and the latter is at most \(H(E\mid W)\).  This proves
\eqref{eq:r-refinement-charge}.
\end{proof}
